\section{Periodization and the exact remainder}\label{sec:remainder}
The approximation error will be expressed through a quotient Duhamel formula. We first establish periodization,
fix the representatives and isometric reductions, and then derive the exact remainder on the quotient.
The analytic estimates needed to bound that remainder are developed in Sections~\ref{sec:analytic} and~\ref{app:o1st}.

On $\Hb$, $-\Delta$ denotes the Friedrichs extension from $C_c^\infty(\Hb)$; on $M$, it denotes the
Friedrichs extension from periodized test functions. The propagator and the powers of the Laplacian on either space
are defined by spectral calculus; we write $\Delta_M$ when the distinction is needed.
The vector field $Y=y\partial_y$ differs from the transport generator $i(y\partial_y-\tfrac12)$ denoted by $Y$ in \cite{S}.

\subsection{Periodization by Cauchy--Schwarz}\label{sec:t2}

\begin{lemma}[uniform Poincar\'e series bound]\label{lem:p1}
For every Fuchsian $\Gamma$ and $s>1$: $\sup_{z}\sum_{\gamma}\ue^{-s\,d(x_0,\gamma z)}<\infty$.
\end{lemma}
\begin{proof}
The orbit $\Gamma x_0$ is uniformly separated, and its stabilizer is finite. Packing balls of a fixed radius
around orbit points gives $\#\{\gamma:d(z,\gamma x_0)<n+1\}\le C\ue^{n}$, uniformly in $z$.
\end{proof}

\begin{lemma}[weighted periodization and transport]\label{lem:t2}
For every Fuchsian $\Gamma$, every $s>1/2$, $t\ge0$, every real phase $\varphi$ and every measurable $w$ with $\ue^{s\tilde\psi}w\in L^2(\Hb)$:
\[
\norm{\bigl[(S(t)w)\,\ue^{i\varphi/\hbar}\bigr]_\Gamma}_{L^2(M)}\le C(\Gamma,s)\,\ue^{st}\,\norm{\ue^{s\tilde\psi}w}_{L^2(\Hb)},
\]
with $C(\Gamma,s)=2^sP_\Gamma(2s)^{1/2}$, where $P_\Gamma(\sigma):=\sup_z\sum_{\gamma\in\Gamma}\ue^{-\sigma d(x_0,\gamma z)}$ is the Poincar\'e series bound of Lemma~\ref{lem:p1}.
\end{lemma}

\begin{proof}

\emph{Pointwise Cauchy--Schwarz.} For $f$ on $\Hb$ and $z\in\Hb$:
\[
\Bigl|\sum_\gamma f(\gamma z)\Bigr|^2\le\Bigl(\sum_\gamma\ue^{-2sd(\gamma z,x_0)}\Bigr)\Bigl(\sum_\gamma\ue^{2sd(\gamma z,x_0)}|f(\gamma z)|^2\Bigr).
\]
The first factor is $\le P_\Gamma(2s)$ for every $z$, by Lemma~\ref{lem:p1} with exponent $2s>1$ (since $d(\gamma z,x_0)=d(z,\gamma^{-1}x_0)$).

\emph{Unfolding.} Integrate over a fundamental domain (Lemma 3.1 of \cite{S}, Tonelli). This gives
\[
\norm{[f]_\Gamma}^2_{L^2(M)}\le P_\Gamma(2s)\int_\Hb\ue^{2sd(z,x_0)}|f(z)|^2\,d\nu(z).
\]
Prove this first for finite partial sums, then pass to the limit by monotone convergence. In particular $\sum_\gamma|f\circ\gamma|$ converges a.e. and in $L^2(M)$. Torsion-freeness is not needed: the fixed sets of elliptic elements are null, so unfolding over a fundamental domain carries no stabilizer factor, and any multiplicity is already counted in the orbit series $P_\Gamma(2s)$.

\emph{Transport.} Take $f=(S(t)w)\ue^{i\varphi/\hbar}$, so $|f|=|S(t)w|$.
Since $(S(t)w)(z)=\ue^{-t/2}w(z_t)$, where $z_t$ is at distance $t$ from $z$, we have
$d(z,x_0)\le d(z_t,x_0)+t$. Unitarity of $S(t)$ therefore gives
\[
  \int\ue^{2sd(z)}|S(t)w(z)|^2\,d\nu(z)
  \le\ue^{2st}\int\ue^{2sd(z')}|w(z')|^2\,d\nu(z').
\]

Finally $\ue^{sd}\le2^s\ue^{s\tilde\psi}$, since $\tilde\psi=\log\cosh d\ge d-\log2$. \qedhere

\end{proof}



\subsection{Regularity and the base point}\label{subsec:dd-base}

The theorems are stated for amplitudes of finite norm and for every boundary point $b$. The next two lemmas reduce them to
$C^\infty$ amplitudes, the base point $i$ and $b=\infty$.

\begin{lemma}[regularity]\label{lem:lift}
Let $0<\alpha\le1$ and $F\in\operatorname{Dom}(\ue^{\alpha\Lambda})$. Then $F$ has a unique $C^\infty$ representative, and every word in
$X,Y$ applied to it lies in $L^2(\Hb)$. Consequently every $a$ with $\norm{\ue^{\alpha\Lambda}X_0^\beta a}<\infty$, and every
$a\in\operatorname{Dom}(\ue^{\alpha\Lambda})$, has a unique $C^\infty$ representative. The approximations $u^{(K)}$ are built from this
representative.
\end{lemma}

\begin{proof}

By Lemma~\ref{lem:wc} at $\alpha''=0$ (with $\tau_0=1$), every word applied to $F$ in the sense of distributions lies in $L^2$. The same holds
for $F\circ p$ for every affine isometry $p(z)=p_1z+p_2$, $p_1>0$: these commute with $X$, $Y$ and $\Lambda$, so
$F\circ p\in\operatorname{Dom}(\ue^{\alpha\Lambda})$.

Let $w\in L^2$ be such that all words applied to $w$ lie in $L^2$, and let $c\in C_c^\infty(\Hb)$ be real. With the Euclidean
Laplacian $\Delta_{\mathrm E}=\partial_x^2+\partial_y^2=y^{-2}(X^2+Y^2-Y)$,
\[
\Delta_{\mathrm E}(cw)=(\Delta_{\mathrm E}c)\,w+2y^{-2}\bigl((Xc)(Xw)+(Yc)(Yw)\bigr)+c\,y^{-2}(X^2+Y^2-Y)w,
\]
a finite sum of terms $c'w'$ of the same kind. So $\Delta_{\mathrm E}^n(cw)\in L^2(\mathbb R^2)$ for every $n$; on $\operatorname{supp}c$
the measures $dx\,dy$ and $d\nu$ are comparable.

Hence $|\xi|^{2n}\widehat{cw}\in L^2(\mathbb R^2)$ for every $n$, all moments of $\widehat{cw}$ are integrable, and $cw$ has a
$C^\infty$ representative.

Take $c=1$ near $i$. This gives a $C^\infty$ representative of $F\circ p$ near $i$ for every $p$, that is, of $F$ near
every point. Two continuous representatives on an open set coincide, so they glue to a $C^\infty$ representative of $F$, which is unique.

For the consequences take $F=X_0^\beta a$, respectively $F=a$; $X_0^{-\beta}$ is smooth and positive. \qedhere

\end{proof}

\begin{lemma}[base point]\label{lem:base}
Let $x_0\in\Hb$ and $b\in\partial\Hb$.
\begin{itemize}
\item[\textup{(i)}] There is $g\in SL(2,\mathbb R)$, unique up to sign, with $g\,x_0=x_0$ and $g\,\infty=b$. We normalize
$\varphi_b:=\varphi_\infty\circ g^{-1}$. The construction of \cite[(1.20)--(1.29)]{S} is covariant under isometries: the initial state and
the approximations for $b$ are $v\mapsto v\circ g^{-1}$ applied to those for $\infty$ with the amplitude $a\circ g$.
\item[\textup{(ii)}] Let $p(z)=(\operatorname{Im}x_0)\,z+\operatorname{Re}x_0$ and $q:=g\,p$, so $q\,i=x_0$ and $q\,\infty=b$. For an amplitude $a$ put
$a_q:=a\circ q$ and $\Gamma_q:=q^{-1}\Gamma q$. Then, with $c:=\ue^{-i\log(\operatorname{Im}x_0)/\hbar}$,
\[
(a\,\ue^{i\varphi_b/\hbar})\circ q=c\,a_q\,\ue^{i\varphi_\infty/\hbar},\qquad u^{(K)}_b(t)[a]\circ q=c\,u^{(K)}_\infty(t)[a_q],
\]
where the subscript names the boundary point of the construction.
\item[\textup{(iii)}] The map $f\mapsto f\circ q$ is unitary from $L^2(\Hb/\Gamma)$ to $L^2(\Hb/\Gamma_q)$, intertwines the periodizations over $\Gamma$ and
$\Gamma_q$, and commutes with $\Delta$, $\cU(t)$ and $\ue^{\alpha\Lambda}$. The reversed-weight and original norms of $a$ with base point $x_0$ equal those of $a_q$
with base point $i$. For every $s$, $\sup_z\sum_{\gamma\in\Gamma_q}\ue^{-s\,d(i,\gamma z)}=\sup_z\sum_{\gamma\in\Gamma}\ue^{-s\,d(x_0,\gamma z)}$.
\end{itemize}
Consequently the error norm of $(\Gamma,x_0,b,a)$ equals that of $(\Gamma_q,i,\infty,a_q)$, and an estimate proved for base point $i$ and
$b=\infty$, with constants depending on the group only through quantities invariant under $\Gamma\mapsto\Gamma_q$, $x_0\mapsto i$, holds
for every $x_0$ and $b$.
\end{lemma}

\begin{proof}
\begin{itemize}
\item (i) The stabilizer of $x_0$ in $PSL(2,\mathbb R)$ acts simply transitively on $\partial\Hb$. The phase $\varphi_b$, the transport and
the operators $\Delta(s)$ for $b$ are defined from the geodesics ending at $b$, so they are the transports by $g$ of those for $\infty$.
\item (ii) $\varphi_b\circ q=\varphi_\infty\circ p=\varphi_\infty-\log\operatorname{Im}x_0$, since $\operatorname{Im}p(z)=(\operatorname{Im}x_0)\operatorname{Im}z$.
By (i), composing with $g$ turns the objects for $b$ and $a$ into those for $\infty$ and $a\circ g$. The affine map $p$ commutes with
$X$, $Y$ and $S(t)$, hence with $\Delta(s)$ and the $P_k(t)$, and it changes $\ue^{i\varphi_\infty/\hbar}$ by the constant factor $c$. This
gives both identities.
\item (iii) $q$ is an isometry of $\Hb$ with $q\,i=x_0$, so $d(x_0,qz)=d(i,z)$ and the weights $X_0$, $\langle d\rangle$ based at $x_0$ pull
back to those based at $i$. The spectral functions of $-\Delta$ commute with composition by isometries. For the last identity,
$d(i,q^{-1}\gamma qz)=d(x_0,\gamma(qz))$. \qedhere
\end{itemize}
\end{proof}

\begin{remark}
The hypothesis $g\,x_0=x_0$ is used in (iii): the Poincar\'e sum of $\Gamma$ is not bounded uniformly in the base point when $\Gamma$
has parabolic elements, so the base point of the weight must be the point at which the sum is taken.
\end{remark}

\subsection{The quotient propagator and Duhamel formula}\label{subsec:dd-surf}

Periodization needs decay, and Proposition~\ref{prop:ws} gives weighted bounds only for $V_\varepsilon$, not for $V$. So $[\ue^{i\varphi/\hbar}V(s)^*g]_\Gamma$ cannot be read as a convergent sum, and the identity $\cU(t)u_\Gamma=\sum_\gamma(\cU(t)u)\circ\gamma^{-1}$ of \cite{S} is not available for $u=\ue^{i\varphi/\hbar}V(s)^*g$. We therefore define this object on $M$ through the surface propagator, and prove the Duhamel formula on $M$.

Let $\omega\ge1$ be a locally bounded measurable weight on $\Hb$, $L^2_\omega:=\{f:\omega f\in L^2(\Hb)\}$, and $\delta_sz:=x+i\ue^sy$. Assume:
\begin{itemize}
\item[(W1)] $f\mapsto f_\Gamma$ is bounded from $L^2_\omega$ to $L^2(M)$, the sum converging absolutely a.e.;
\item[(W2)] $\sup\{\omega(z)/\omega(\delta_sz):|s|\le T,\ z\in\Hb\}<\infty$ for every $T$.
\end{itemize}
Lemma~\ref{lem:t2} (its proof, $t=0$) gives (W1) for $\omega=W_\sigma$, $\sigma>1/2$. Proposition~\ref{os:p1} gives it for the periodization weight $\omega=X_0^{1/2}\ell^{\,c'}$, $\ell=\log(2X_0)/\log2$, of Section~\ref{sec:printed}. Since $d(z,\delta_sz)=|s|$ and $\tilde\psi$ is $1$-Lipschitz, (W2) holds with the explicit ratios $W_\sigma(z)/W_\sigma(\delta_sz)\le\ue^{\sigma|s|}$ and $\omega(z)/\omega(\delta_sz)\le\ue^{|s|/2}(1+|s|/\log 2)^{c'}$. By (W2), $S(s)$ is bounded on $L^2_\omega$, locally uniformly in $s$; it is strongly continuous there, by density of $C_c(\Hb)$.

Let $\mathcal D_\omega$ be the set of $g\in L^2_\omega$ with $Xg,Yg,X^2g,Y^2g\in L^2_\omega$ (weak derivatives). For $g\in L^2_\omega$ and $s\in\mathbb R$ put
\[
\mathcal V_sg:=\ue^{-is/(2\hbar)}\,\cU(-s)\,\bigl[\ue^{i\varphi/\hbar}S(s)g\bigr]_\Gamma\in L^2(M),
\]
with $\cU$ the propagator on $M$. Thus
\[
\mathcal V_0g=[\ue^{i\varphi/\hbar}g]_\Gamma,\qquad
\norm{\mathcal V_sg}_{L^2(M)}=\norm{[\ue^{i\varphi/\hbar}S(s)g]_\Gamma}_{L^2(M)}.
\]
This $\mathcal V_sg$ is what $[\ue^{i\varphi/\hbar}V(s)^*g]_\Gamma$ formally means by \cite[(1.23)]{S}.

Put $\Delta_\varphi:=\ue^{-i\varphi/\hbar}\Delta\ue^{i\varphi/\hbar}$ and
$\Delta_\varphi(s):=S(s)^{-1}\Delta_\varphi S(s)$, so that
\[
  \Delta_\varphi(s)=\Delta(s)-\frac{2i}{\hbar}(Y-\tfrac12)-\hbar^{-2}.
\]

\begin{lemma}[Laplacian on the weighted class]\label{lem:surf-lap}
For $g\in\mathcal D_\omega$ and $s\in\mathbb R$, $[\ue^{i\varphi/\hbar}S(s)g]_\Gamma\in\operatorname{Dom}(\Delta_M)$ and $\Delta_M[\ue^{i\varphi/\hbar}S(s)g]_\Gamma=[\ue^{i\varphi/\hbar}S(s)\Delta_\varphi(s)g]_\Gamma$. Hence $\mathcal V_sg\in\operatorname{Dom}(\Delta_M)$ and $\Delta_M\mathcal V_sg=\mathcal V_s\Delta_\varphi(s)g$.
\end{lemma}

\begin{proof}

Put $f:=\ue^{i\varphi/\hbar}S(s)g$. Then
\[
\begin{aligned}
Xf&=\ue^{-s}\ue^{i\varphi/\hbar}S(s)Xg,\\
Yf&=\ue^{i\varphi/\hbar}S(s)(Y-i/\hbar)g,\\
\Delta f&=\ue^{i\varphi/\hbar}S(s)\Delta_\varphi(s)g.
\end{aligned}
\]
So $f,Xf,Yf,\Delta f\in L^2_\omega$ by (W2), and $f\in\operatorname{Dom}(\Delta)$ because $\Delta$ is essentially self-adjoint on $C_c^\infty(\Hb)$.

Take radial $\chi_n=\rho(d(x_0,\cdot)/n)$, with $\rho\in C_c^\infty$, $\rho=1$ near $0$. Then $0\le\chi_n\le1$, $\chi_n\to1$, and $|X\chi_n|+|Y\chi_n|+|\Delta\chi_n|\le C/n$. Since
\[\Delta(\chi_nf)=\chi_n\Delta f+2(X\chi_n\,Xf+Y\chi_n\,Yf)+(\Delta\chi_n)f,\]
dominated convergence gives $\chi_nf\to f$ and $\Delta(\chi_nf)\to\Delta f$ in $L^2_\omega$.

Fix $n$. Essential self-adjointness gives test functions $F_m\to\chi_nf$ in graph norm. With a fixed $\theta\in C_c^\infty(\Hb)$, $\theta=1$ on $\operatorname{supp}\chi_n$, the test functions $\theta F_m$ converge to $\chi_nf$ in graph norm, since $\norm{\nabla F}^2=\langle F,-\Delta F\rangle$ controls the cross term. Their supports lie in the compact set $\operatorname{supp}\theta$, where $\omega$ is bounded.

For a test function $F$, $F_\Gamma\in\operatorname{Dom}(\Delta_M)$ and $\Delta_MF_\Gamma=(\Delta F)_\Gamma$. By (W1) and closedness of $\Delta_M$, first in $m$ and then in $n$, the claim for $f$ follows. The statement for $\mathcal V_s$ follows since $\cU(-s)$ commutes with $\Delta_M$.\qedhere

\end{proof}

\begin{lemma}[Duhamel on $M$]\label{lem:surf-duh}
Let $g\in C^1([0,T];L^2_\omega)$ with $g(s)\in\mathcal D_\omega$ and $Yg,X^2g,Y^2g\in C([0,T];L^2_\omega)$. Then, with $c=i\hbar/2$,
\[
\mathcal V_Tg(T)-\mathcal V_0g(0)=\int_0^T\mathcal V_s\bigl(g'(s)-c\,\Delta(s)g(s)\bigr)\,ds\qquad\text{in }L^2(M).
\]
\end{lemma}

\begin{proof}

For $g_0$ with $g_0,Yg_0\in L^2_\omega$: $S(s)g_0-S(r)g_0=\int_r^sS(\tau)(Y-\tfrac12)g_0\,d\tau$. This holds in $L^2(\Hb)$ (generator of the dilation group), the right side converges in $L^2_\omega$ by strong continuity, and $L^2_\omega\subset L^2$ since $\omega\ge1$.

Hence $s\mapsto S(s)g(s)$ is $C^1$ into $L^2_\omega$, and by (W1) $h(s):=[\ue^{i\varphi/\hbar}S(s)g(s)]_\Gamma$ is $C^1$ into $L^2(M)$ with $h'=[\ue^{i\varphi/\hbar}S(s)((Y-\tfrac12)g+g')]_\Gamma$.

By Lemma~\ref{lem:surf-lap}, $h(s)\in\operatorname{Dom}(\Delta_M)$ and $\Delta_Mh=[\ue^{i\varphi/\hbar}S(s)\Delta_\varphi(s)g]_\Gamma$, continuous in $s$. So $\tfrac{d}{ds}\cU(-s)h(s)=\cU(-s)\bigl(h'(s)-\tfrac{i\hbar}{2}\Delta_Mh(s)\bigr)$.

Since $\Delta_\varphi(s)=\Delta(s)-\tfrac{2i}{\hbar}(Y-\tfrac12)-\hbar^{-2}$, we have $-\tfrac{i\hbar}{2}\Delta_\varphi(s)=-c\Delta(s)-(Y-\tfrac12)+\tfrac{i}{2\hbar}$. The $(Y-\tfrac12)$ terms cancel, and $\tfrac{i}{2\hbar}$ cancels the derivative of $\ue^{-is/(2\hbar)}$. Thus $\tfrac{d}{ds}\mathcal V_sg(s)=\mathcal V_s(g'(s)-c\Delta(s)g(s))$; integrate.\qedhere

\end{proof}

For $k\ge1$ put $\Sigma_k(t):=\{t>s_1>\cdots>s_k>0\}$ and let $c:=i\hbar/2$.
\begin{corollary}[Dyson expansion on $M$]\label{cor:surf-dyson}
If every word in $X,Y$ applied to $a$ lies in $L^2_\omega$, then with $a^{(K)}(t)=\sum_{k\le K}c^kP_k(t)a$,
\[
\mathcal V_0a-\mathcal V_ta^{(K)}(t)=c^{K+1}\int_{\Sigma_{K+1}(t)}\mathcal V_{s_1}\bigl(\Delta(s_1)\cdots\Delta(s_{K+1})a\bigr)\,ds.
\]
\end{corollary}

\begin{proof}
$P_k(s)a$ is a finite sum of continuous scalar functions times fixed words applied to $a$, so Lemma~\ref{lem:surf-duh} applies to it. Since $\partial_sP_{k}(s)=\Delta(s)P_{k-1}(s)$ for $k\ge1$ and $P_k(0)=0$, with $J_k(t):=\int_0^t\mathcal V_s\Delta(s)P_k(s)a\,ds$ we get $\mathcal V_tP_k(t)a=J_{k-1}(t)-cJ_k(t)$ for $k\ge1$, and $\mathcal V_ta-\mathcal V_0a=-cJ_0(t)$. Multiplying by $c^k$ and summing telescopes to $\mathcal V_0a-\mathcal V_ta^{(K)}(t)=c^{K+1}J_K(t)$. Expanding $P_K$ gives the integral over $\Sigma_{K+1}(t)$.
\end{proof}


\subsection{Dyson identities and the exact remainder}\label{subsec:dd-id}

Let $V(t)$ be the unitary operator of \eqref{eq:factor}, let $V(t,s):=V(t)V(s)^{-1}$, the propagator of
$i\partial_tV=-\tfrac{\hbar}{2}\Delta(t)V$ (\cite[(1.22)]{S}), and let $c:=i\hbar/2$. By \eqref{eq:factor} (which is \cite[(1.23)]{S}),
\[
\cU(t)(a\,\ue^{i\varphi/\hbar})=\ue^{-it/(2\hbar)}\bigl(S(t)V(t)a\bigr)\ue^{i\varphi/\hbar},
\qquad (S(t)w)(x,y)=\ue^{-t/2}w(x,\ue^{t}y),
\]
with $S(t)$ unitary. Then
\[S(t)^{-1}XS(t)=\ue^{-t}X,\qquad S(t)^{-1}YS(t)=Y,\]
hence $S(t)^{-1}\Delta S(t)=\Delta(t)$.
Multiplication by $\ue^{i\varphi/\hbar}=y^{-i/\hbar}$ preserves every Sobolev space $H^k$, and so do $\cU(t)$ and $S(t)$.
By the factorization, so do $V(t)^{\pm1}$ and $V(t,s)$.

\begin{proposition}[Dyson identities and quotient remainder]\label{prop:dd-id}
\begin{enumerate}
\item[\textup{(i)}] \emph{Dyson expansion.} Let $a$ be such that every word in $X,Y$ applied to $a$ lies in $L^2(\Hb)$, and put
\[
P_k(t)=\int_{t>s_1>\cdots>s_k>0}\Delta(s_1)\cdots\Delta(s_k)\,ds,\qquad a^{(K)}(t)=\sum_{k=0}^{K}c^kP_k(t)a .
\]
Then, with $\Sigma_{K+1}(t)=\{t>s_1>\cdots>s_{K+1}>0\}$,
\[
a-V(t)^*a^{(K)}(t)=c^{K+1}\int_{\Sigma_{K+1}(t)}V(s_1)^*\,\Delta(s_1)\cdots\Delta(s_{K+1})a\,ds .
\]
This $P_k$ is exactly the $P_k$ of \cite{S}
(the display before the semiclassical approximation of order $K$). Every amplitude $a$ of Theorem~\ref{thm:dd-appx} satisfies the
hypothesis, by Lemma~\ref{lem:sc}.
\item[\textup{(ii)}] \emph{Transport of $\Delta^N$} (the conjugated operator is that of \cite[Proposition~2.2 and the proof of Proposition~3.10, (3.31)--(3.34)]{S}). Put $\Delta_\varphi:=e^{-i\varphi/\hbar}\Delta e^{i\varphi/\hbar}$, so that
$\hbar^2\Delta_\varphi=\hbar^2\Delta-2i\hbar Y-1+i\hbar$. Put $\Delta_\varphi(s):=S(s)^{-1}\Delta_\varphi S(s)$, so that
$\hbar^2\Delta_\varphi(s)=\hbar^2\Delta(s)-2i\hbar Y-1+i\hbar$. Then for every $b$ such that every word in $X,Y$ applied to $b$ lies in $L^2(\Hb)$,
\[
\Delta_\varphi^N V(s)^*b=V(s)^*\Delta_\varphi(s)^Nb.
\]
\item[\textup{(iii)}] \emph{The error.} Let $E:=\Delta^N\bigl[\cU(t)(a\,\ue^{i\varphi/\hbar})_\Gamma-u^{(K)}(t)_\Gamma\bigr]$, with $u^{(K)}$ as in \cite{S}. If every word in $X,Y$ applied to $a$ lies in $L^2_\omega$, for a weight $\omega$ as in Section~\ref{subsec:dd-surf}, then
\[
\begin{gathered}
\norm{E}_{L^2(M)}\le\Bigl(\frac{\hbar}{2}\Bigr)^{K+1}\int_{\Sigma_{K+1}(t)}\bigl\|\mathcal V_{s_1}b'(s)\bigr\|_{L^2(M)}\,ds,\\
b'(s):=\Delta_\varphi(s_1)^N\Delta(s_1)\cdots\Delta(s_{K+1})a.
\end{gathered}
\]
Everything of high order is now inside $b'(s)$, which is a fixed word-polynomial in $X,Y$ applied to $a$.
\end{enumerate}
\end{proposition}

\begin{proof}
\medskip\noindent\textbf{(i).}

$V(s)^*=V(0,s)$ and $\partial_sV(t,s)=-c\,V(t,s)\Delta(s)$ on $H^\infty$. So for $b$ with all words in $L^2$, $s\mapsto V(s)^*b$ is
differentiable in $L^2$ with derivative $-c\,V(s)^*\Delta(s)b$. (In integrated form this is Lemma~\ref{lem:surf-duh} for the trivial
group, the weight $\omega=1$ and the constant path $g(s)=b$: then $\mathcal V_sg=\ue^{i\varphi/\hbar}V(s)^*g$ by \eqref{eq:factor}; that lemma was proved directly on the quotient.)

$P_k(s)a$ is a finite sum of continuous scalar functions of $s$ times fixed words applied to $a$; $\partial_sP_k(s)=\Delta(s)P_{k-1}(s)$
and $P_k(0)=0$ for $k\ge1$.

With $J_k(t):=\int_0^tV(s)^*\Delta(s)P_k(s)a\,ds$ this gives $V(t)^*P_k(t)a=J_{k-1}(t)-cJ_k(t)$ for $k\ge1$, and $V(t)^*a-a=-cJ_0(t)$.

Multiplying by $c^k$ and summing telescopes to $a-V(t)^*a^{(K)}(t)=c^{K+1}J_K(t)$. Expanding $P_K$ gives the integral over $\Sigma_{K+1}(t)$.

\medskip\noindent\textbf{(ii).}

For $c'\in H^2$, the commutation of the propagator with the Laplacian gives the following identity.

$\cU(s)$ commutes with $\Delta$, so $\cU(s)(e^{i\varphi/\hbar}\Delta_\varphi c')=\Delta\,\cU(s)(e^{i\varphi/\hbar}c')$.

By \cite[(1.23)]{S} this reads $S(s)V(s)\Delta_\varphi c'=\Delta_\varphi S(s)V(s)c'$, i.e. $V(s)\Delta_\varphi=\Delta_\varphi(s)V(s)$ on $H^2$.

Multiplying by $V(s)^*$ on both sides is legitimate because $V(s)^*$ preserves $H^2$. Iterating $N$ times gives the assertion, since $b$ lies in every $H^k$.

\medskip\noindent\textbf{(iii).}

By definition $u^{(K)}(t)=\ue^{-it/(2\hbar)}\ue^{i\varphi/\hbar}S(t)a^{(K)}(t)$, so $u^{(K)}(t)_\Gamma=\cU(t)\mathcal V_ta^{(K)}(t)$. Since $\cU(t)$ is unitary on $L^2(M)$ and commutes with $\Delta_M$, $\norm E=\norm{\Delta_M^N(\mathcal V_0a-\mathcal V_ta^{(K)}(t))}$. (No commutation of $\cU(t)$ with periodization is used.)

By Corollary~\ref{cor:surf-dyson}, $\mathcal V_0a-\mathcal V_ta^{(K)}(t)=c^{K+1}\int_{\Sigma_{K+1}(t)}\mathcal V_{s_1}(\Delta(s_1)\cdots\Delta(s_{K+1})a)\,ds$.

By Lemma~\ref{lem:surf-lap} $N$ times, $\Delta_M^N\mathcal V_{s_1}g=\mathcal V_{s_1}\Delta_\varphi(s_1)^Ng$. Closedness of $\Delta_M^N$ and Minkowski's inequality give the bound, with $|c|^{K+1}=(\hbar/2)^{K+1}$.
(i) and (ii) are identities on $\Hb$; on $M$ the proof of (iii) uses their surface versions in Section~\ref{subsec:dd-surf}. \qedhere
\end{proof}

\subsection{Local pointwise control}
\begin{lemma}[local Sobolev inequality on a Fuchsian quotient]\label{pc:g43}
Let $m(z):=\#\{\gamma\in\Gamma:\ d(\tilde z,\gamma\tilde z)<2\}$ for $z\in M$ with a lift $\tilde z$; $m$ is finite and does not depend on the lift.
There is an absolute constant $C$ such that every $u\in\operatorname{Dom}(\Delta_M)$ has a continuous representative, and for all $z\in M$,
\[
  |u(z)|\le C\,m(z)^{1/2}\big(\norm{\Delta u}_{L^2(M)}+\norm{u}_{L^2(M)}\big).
\]
\end{lemma}
\begin{proof}
\emph{Step 1 (the lift).} Let $\tilde u$ and $\tilde f$ be the $\Gamma$-invariant lifts of $u$ and of $f:=\Delta_Mu$. Then $\Delta\tilde u=\tilde f$ in the
sense of distributions on $\Hb$: for $\phi\in C_c^\infty(\Hb)$, unfolding gives
$\langle\tilde u,\Delta\phi\rangle_{\Hb}=\langle u,(\Delta\phi)_\Gamma\rangle_M=\langle u,\Delta_M\phi_\Gamma\rangle_M=\langle f,\phi_\Gamma\rangle_M=\langle\tilde f,\phi\rangle_{\Hb}$.

\emph{Step 2 (counting).} For a $\Gamma$-invariant $\tilde h$ and a fundamental domain $F$,
$\int_{B(\tilde z,1)}|\tilde h|^2d\nu=\int_Fn(w)|h(w)|^2d\nu(w)$, where $n(w):=\#\{\gamma:\gamma w\in B(\tilde z,1)\}$. If $\gamma_0w$ and $\gamma w$ both lie in
$B(\tilde z,1)$, then $d(\tilde z,\gamma\gamma_0^{-1}\tilde z)<2$; hence $n(w)\le m(z)$ and $\int_{B(\tilde z,1)}|\tilde h|^2\le m(z)\norm{h}^2_{L^2(M)}$. In
particular $\tilde u$ and $\tilde f$ are locally square integrable. The number $m(z)$ is finite because $\Gamma$ acts properly discontinuously.

\emph{Step 3 (interior regularity).} Let $v\in L^2(B(\tilde z,1))$ with $\Delta v\in L^2(B(\tilde z,1))$ in the sense of distributions. By interior
elliptic regularity, uniform in $\tilde z$ by the homogeneity of $\Hb$, $v\in H^2(B(\tilde z,\tfrac12))$ with
$\norm v_{H^2(B(\tilde z,1/2))}\le C(\norm{\Delta v}_{L^2(B(\tilde z,1))}+\norm v_{L^2(B(\tilde z,1))})$, and since $H^2\subset C^0$ in dimension two, $v$ is
continuous on $B(\tilde z,\tfrac12)$ with $|v(\tilde z)|\le C\norm v_{H^2(B(\tilde z,1/2))}$.

Apply Step~3 to $v=\tilde u$ and use Step~2 for $\tilde h=\tilde u$ and $\tilde h=\tilde f$. The continuous representative of $\tilde u$ is $\Gamma$-invariant
and descends to $M$. \qedhere
\end{proof}
